\label{sec:history}
\subsection{The direction that failed}
For thirty years the main line of attack sought an \emph{invariant that counts generators}. Candidates were Voiculescu's free entropy
dimensions ($\delta,\delta_0,\delta^*,\delta^\star$), Connes--Shlyakhtenko $L^2$-Betti numbers of the algebra, and cost-like quantities. On canonical generators they
all return $n$ (or $n-1$). The expert's remark, that ``we have been heading in the wrong direction'', reads naturally against this background. The paper
shows that the count is a property of the \emph{generating set}, not of the algebra:
\begin{itemize}[leftmargin=*]
\item a free Haar $n$-tuple, with $\delta_0=n$, generates $L(\F_{n+1})$;
\item $L(\F_2)$ has finite generating tuples with $\delta=\delta_0=\delta^*=\delta^\star=n$ for every $n\ge2$ (Corollary 7.1).
\end{itemize}
The resolution came instead from \emph{dynamics}: an explicit, rich family of automorphisms and an adaptive limit. That is a constructive object, not a
numerical invariant.

\subsection{What stays true, and what must change}
\begin{itemize}[leftmargin=*]
\item \emph{Upper-bound theorems survive.} These are results that bound $\delta_0$ for \emph{all} generating sets, such as Voiculescu's absence of Cartan
subalgebras, Ge's primeness, Jung's strong $1$-boundedness, and Hayes' $1$-bounded entropy. They are statements about every generating set, and remain
valid. $L(\F_n)$ still has no Cartan subalgebra, is prime and strongly solid, and $h(L(\F_n))=\infty$ for all $n$.
\item \emph{Generator-counting invariants must collapse.} Any quantity that is an invariant of the von Neumann algebra and was expected to equal
$n$ or $n-1$ on $L(\F_n)$ must take a single value for all $n$. For example, the Connes--Shlyakhtenko first $L^2$-Betti number of $L(\F_n)$ cannot
equal $n-1$ for every $n$.
\item \emph{Fundamental group.} $\mathcal F(L(\F_2))=\R_{>0}$; every interpolated factor is isomorphic to every other, including $L(\F_\infty)$.
\item \emph{Generation is not stable under small perturbations.} Choosing $r_0$ small in the construction gives free Haar $n$-tuples \emph{arbitrarily close in
operator norm} to the canonical generators $\lambda(x_1),\dots,\lambda(x_n)$ of $L(\F_{n+1})$ that generate all of it. The canonical tuple generates a proper
subfactor with a free complement. Generation is invisible to any norm-local count.
\end{itemize}

\subsection{Relation to earlier techniques}
\begin{itemize}[leftmargin=*]
\item \emph{Voiculescu's cyclomorphy} exponentiates trace-preserving polynomial vector fields on free semicircular generators. The paper redoes it for
Haar unitaries with a fixed bounded logarithm, proving descent through the unitary relations from moments.
\item \emph{Guionnet--Shlyakhtenko free monotone transport} builds isomorphisms by changes of variables at a \emph{fixed} rank. Here the rank changes
through an infinite product of small automorphisms and a limiting generation argument.
\item \emph{Shlyakhtenko's Fuchsian preprint} (September 2026) finds free complements by finite automorphisms plus limiting generation. The idea was
in the air, and the small-cocycle absorption is the new ingredient.
\item \emph{Fox calculus} (1953) supplies the cocycle rule, and \emph{Ricard--Xu} the free Khintchine decomposition.
\end{itemize}
